% Generated draft: requires mathematical and source review.
\documentclass[11pt]{article}
\usepackage[margin=25mm]{geometry}
\usepackage{fontspec}
\setmainfont{Latin Modern Roman}
\usepackage{amsmath,amssymb,mathtools,unicode-math}
\setmathfont{Latin Modern Math}
\usepackage{enumitem,needspace,xurl,hyperref}
\setlength{\parindent}{0pt}
\setlength{\parskip}{5pt}
\setlength{\emergencystretch}{3em}
\newcommand{\N}{\mathbb{N}}
\newcommand{\Z}{\mathbb{Z}}
\newcommand{\Q}{\mathbb{Q}}
\newcommand{\R}{\mathbb{R}}
\newcommand{\C}{\mathbb{C}}
\newcommand{\F}{\mathbb{F}}
\newcommand{\A}{\mathbb{A}}
\newcommand{\PP}{\mathbb{P}}
\newcommand{\E}{\mathbb{E}}
\newcommand{\eps}{\varepsilon}
\newcommand{\dd}{\,\mathrm{d}}
\begin{document}
\section*{UnsolvedMath ID: 8800083}
\textbf{Source problem code:} catalogue entry
\subsection{Definitions and mathematical statement}
For a group $G$, element inversion is the map $g\mapsto g^{-1}$ characterized by $gg^{-1}=g^{-1}g=1$. The source asks for groups in which computing this inverse is infeasible under reasonable cryptographic assumptions. An abstract group alone has no computational complexity; element encodings and allowed operations must be part of the construction.

A possible cryptographic framework would use a security-parameterized family of encoded groups, efficient multiplication and identity representation, and a specified sampler for challenge elements. The goal would then be hardness of producing an encoding of $g^{-1}$ from an encoding of a sampled $g$, under an explicit assumption. The source does not prescribe this exact framework, sampling distribution, success criterion, or classical/quantum adversary model; nor does it state whether the groups are finite or whether trapdoors are allowed.

A relevant model constraint is that, if a suitable positive multiple $M$ of an element's order is efficiently known, then
\[
 g^{-1}=g^{M-1},
\]
which can be computed by repeated squaring in a number of multiplications polynomial in $\log M$. Thus an order/representation convention is essential to any proposed hard-inversion construction; no universal impossibility conclusion is asserted. Hardness of discrete logarithms, undoing a group action, or inverting a different one-way function is not silently substituted for element inversion.
\subsection{Short English statement}
Construct an effectively usable encoded group whose element-inversion problem is provably hard from a stated cryptographic assumption. The encoding, operations, input distribution, and hardness notion are missing from the source.

The supplied summary reports no verified construction and says common group representations have easy inversion. It mentions infeasible-inversion and one-way-action themes without an exact theorem; those notions are not equated here. The OPEN-TRIAGE label and 2006 workshop question are retained with current primary-literature audit required.
\subsection{Sources}
\begin{itemize}
\item Local UnsolvedMath registry, permanent ID 8800083, source code catalogue entry, ``Groups with infeasible inversion,'' proposed by Ronald Rivest. Dataset curation and metadata: Ulam AI and UnsolvedMath Contributors, CC BY 4.0.
\item IPAM, \emph{Workshop on Number Theory and Cryptography --- Open Problems} (2006), Attachment 6243, legacy PPT logical slide 17, schedule locator \url{https://www.ipam.ucla.edu/programs/workshops/workshop-i-number-theory-and-cryptography-open-problems/?tab=schedule}; pdf-and-legacy-ppt visual-audit extraction, accessed 2026-07-29.
\item The supplied summary mentions Rivest and collaborators, ``signed/unsigned'' ideas, and one-way group actions without a complete primary citation or exact infeasible-inversion result.
\item Source provenance reports inspection of 17 PDFs and two legacy PPT decks among 19 surviving attachments, excluding repetitions and solved definitions; this is supplied provenance, not a new inspection claimed here.
\item Current primary-literature status audit and redistribution terms NEEDS\_REVIEW for the publicly accessible author/publisher copy, separately from the dataset licence.
\end{itemize}
\end{document}
